\subsection{Proposition 2.7.3 of the master prompt: memorization and sideways generalization}
\label{proof:P9}
The statement ``a memorized behaviour has no sideways generalization'' is
false for memorization as defined (Definition~\ref{def:memo}); it is true for a
stronger, leaf-level form.

\textbf{Proposition P9.} (a) \emph{Counterexample.} For every $\tau>0$ there
is a system and a behaviour $b$ that is memorized at $\tau$, has a non-empty
set of sideways inputs, and answers every sideways input correctly.
(b) \emph{Leaf-level memorization.} Say $b$ is \emph{leaf-memorized} if every
non-input node $C$ of $\G$ that is active on some training input of $b$ has an
empty training range, $\mathrm{Rg}(C)=\emptyset$ (no value of $C$ is taken on
$n_{\min}$ or more distinct training inputs). Then no new input is lower or
sideways for $b$: every new input whose derivation contains a non-input node
is diagonal, and so is every new input whose output is not mediated by the
graph.

\textbf{Proof.} (a) Let $m>1/\tau$ behaviours $b_1,\dots,b_m$ with equal weight
$1/m<\tau$. Behaviour $b_1$ computes $C=(A+B)\bmod q$ through a private node $C$
used by no other behaviour, from nodes $A,B$ computed from disjoint inputs,
and reads $y=C$; behaviours $b_2,\dots,b_m$ read other inputs through other
private nodes. Training withholds a set $W$ of $(A,B)$ combinations while
every value of $A$ and of $B$ is seen ($q\ge3$, $|W|$ small). The block that
computes $C$ is the full addition table (the compositional variant of planted
system V3). Every node on the derivation of $b_1$ ($A$, $B$, $C$) is used only
by $b_1$, so $g\le 1/m<\tau$, $\mu(x)<\tau$ on every input, and $b_1$ is
memorized at $\tau$. An input with $(A,B)\in W$ has all node values in range
and the parent tuple of $C$ unseen, so it is sideways; the system writes
$(A+B)\bmod q$ there, which is correct.
(b) Let $x^\ast$ be new. If $D_b(x^\ast)$ has a non-input node $C$, then
$C(x^\ast)$ is a value of $C$ and $\mathrm{Rg}(C)=\emptyset$, so
$C(x^\ast)\notin\mathrm{Rg}(C)$; by Definition~\ref{def:moves} $x^\ast$ is
neither lower nor sideways, hence diagonal (operationally, the out-of-range
value is reported). If $D_b(x^\ast)$ has no non-input node, the output is not
mediated by the graph, which is diagonal by the operational rule of
Definition~\ref{def:moves}(iii) (graph mediation share $0<1-\varepsilon_{\mathrm{cov}}$).
\hfill$\square$

\textbf{Consequences.} The structural memorization verdict is a statement
about reuse across the behaviour suite, not about compositionality inside one
behaviour: (a) shows it can label a compositional but unshared computation as
memorized. That is a design property of the definition, and it fixes what the
K4/K9 memorization claims may say: a model judged memorized is predicted to
have no \emph{shared} principles, not to fail on every sideways input. The
leaf-level form (b) is the one that predicts failure, through
Conjecture~\ref{conj:diag}.

\textbf{Measurable prediction.} (i) On systems trained with a planted
memorized subset (V2 and its trained variant), inputs whose derivation is
leaf-memorized are classified diagonal by $r.\mathrm{move}$ when perturbed to
unseen neighbours, and fail. (ii) The counterexample family exists in real
suites: behaviours judged memorized at $\tau$ whose private nodes pass the
interchange test on unseen tuples generalize sideways; K3 reports how often.
